
\begin{proposition}[The ruling genus is at most two]\label{cdp:st:C.2.CR.2}
In the situation of \ref{cdp:st:C.2.CR.1}, the ruling genus is $g=1$ or
$g=2$. Let $p=\pi(C)$ and $q=\eta(p)$. When $g=2$, the
normalization has just the one point $p$ over $q$, and the
full conductor-image scheme $\Gamma$ is smooth and reduced
at $q$. Its component through $q$ has generic transversal
either an ordinary node or an ordinary cusp.
\end{proposition}
\begin{proof}
Work over a small neighbourhood of $q$ and fix a generator
of the invertible dualizing sheaf $\omega_S$. Finite
duality then identifies the conductor on the normal branch
$(T,p)$ with its canonical module:
\[
 I_p=\omega_{T,p}\subset\OO_{T,p}.
\]
The conductor ideal on $S$ is the direct sum of these
normalization-branch ideals; thus $I_p$ itself embeds
as an ideal of $R=\OO_{S,q}$ by putting zero on the other
normalization branches.

Put $J=(\pi_*\omega_Z)_p$. Proper duality and
Grauert--Riemenschneider vanishing give
\[
 \operatorname{length}(I_p/J)
   =\operatorname{length}(R^1\pi_*\OO_Z)_p=g.
 \tag{C.CR.10}\label{cdp:eq:C.CR.10}
\]
These are the analytic duality and vanishing maps described
in \cite[Sections~2.2.2--2.2.5]{vanStraten1987}.
For precision, the defect of $\OO_T\to R\pi_*\OO_Z$ is
the finite-length module $R^1\pi_*\OO_Z$ in degree one;
duality turns it into the quotient
$\omega_T/\pi_*\omega_Z$, since the positive direct
images of $\omega_Z$ vanish. Finite-length duality
preserves length.

The residue sequence for the smooth exceptional section is
\[
 0\longrightarrow\omega_Z\longrightarrow\omega_Z(C)
   \longrightarrow\omega_C\longrightarrow0 .
 \tag{C.CR.11}\label{cdp:eq:C.CR.11}
\]
The same vanishing theorem makes its pushforward quotient
at $p$ equal to $H^0(C,\omega_C)$, of dimension $g$.
The sheaf $\pi_*\omega_Z(C)$ embeds in $\omega_T$:
away from the exceptional set its sections are holomorphic
two-forms, and the canonical module of a normal surface
extends these across its isolated singular points.
Comparing lengths with \eqref{cdp:eq:C.CR.10} proves
\[
 I_p=(\pi_*\omega_Z(C))_p,\qquad
 I_p/J\simeq H^0(C,\omega_C).
 \tag{C.CR.12}\label{cdp:eq:C.CR.12}
\]
In the fixed dualizing trivialization the exact comparison
of \ref{cdp:st:A.5.3} reads
$\omega_Z=\OO_Z(-2C-V)$, so that
$\omega_Z(C)=\OO_Z(-C-V)$. The quotient in \eqref{cdp:eq:C.CR.12}
is consequently the leading coefficient of order one
along $C$, with the vertical vanishing conditions retained.

Let $\mathfrak m$ be the maximal ideal of $R$. Every
element of $\mathfrak m$ pulls back to a function vanishing
on $C$, because $h(C)=q$. Its square therefore vanishes
to order at least two on $C$. Formula \eqref{cdp:eq:C.CR.12}, including
all vertical conditions in $V$, implies
\[
                 I_p\cap\mathfrak m^2\subset J.
 \tag{C.CR.13}\label{cdp:eq:C.CR.13}
\]
Indeed an element of $I_p$ is already a section of
$\OO_Z(-C-V)$; an additional zero along $C$ puts it
in $\OO_Z(-2C-V)$.

It follows that
\[
 g=\dim(I_p/J)
 \le \dim\bigl(I_p/(I_p\cap\mathfrak m^2)\bigr)
 \le \dim(\mathfrak m/\mathfrak m^2)\le3.
 \tag{C.CR.14}\label{cdp:eq:C.CR.14}
\]
If the middle dimension were three, the ideal $I_p$
would generate the whole cotangent space. Nakayama's
lemma applied to $\mathfrak m/I_p$ would give
$I_p=\mathfrak m$. The full conductor ideal, which
contains $I_p$, would then be $\mathfrak m$, making its
image zero-dimensional. This contradicts the pure
one-dimensional Cohen--Macaulay conductor image.
Thus the middle dimension is at most two, proving
$g\le2$.

When $g=2$, choose $x,y\in I_p$ whose linear parts
are independent, and complete them to coordinates
$(x,y,z)$ on the smooth ambient threefold. The full
conductor ideal contains $(x,y)$ in $R$. The defining
equation of $S$ belongs to $(x,y)$, since the conductor
has dimension one. The quotient by $(x,y)$ is therefore
$\mathbb C\{z\}$. A strictly larger proper conductor
ideal would give a zero-dimensional quotient of this
ring. Purity hence forces the full conductor ideal
to be exactly $(x,y)$, proving smoothness and
reducedness of $\Gamma$ at $q$.

There is no other normalization point over $q$.
The two elements $x,y$ were chosen in the single
summand $I_p$, so both vanish identically on any other
local normalization branch. Such a two-dimensional
hypersurface branch cannot be contained in the
one-dimensional smooth germ $\{x=y=0\}$.

Finally conductor duality gives
$Q\simeq\omega_\Gamma\otimes\omega_S^{-1}|_\Gamma$.
Near the smooth reduced curve $\Gamma$ this is a
line bundle, so the transverse normalization quotient
has length one. A reduced plane curve of delta
invariant one is a node or a cusp. One can see this
directly: multiple branches require exactly two
smooth branches with intersection multiplicity one;
a single branch has normalization
$\mathbb C\{u\}$ and a colength-one subring, necessarily
$\mathbb C\{u^2,u^3\}$.
\end{proof}


